% =========================================================================
% PLATFORM TAXONOMY ROUTING (patho.dentnet.space)
% =========================================================================
% DOC_ID: traumatic_lesions
% DOC_SLUG: traumatic_lesions
% DOC_CATEGORY: cat_traumatic
% DOC_TITLE: الگوریتم جامع: ضایعات تروماتیک و انسدادی بزاقی دهان
% DOC_ALIASES: traumatic lesions, موکوسل, رانولا, سیالولیتیازیس, نوروم تروماتیک, غدد بزاقی, sialolithiasis, mucocele, ranula
% =========================================================================
% === CLINICAL EVIDENCE LEDGER & PMIDS ===
% 1. WHO Classification of Head and Neck Tumours (5th Ed, 2024) - Non-Neoplastic Salivary & Traumatic Lesions.
% 2. Neville BW, Damm DD, Allen CM, Chi AC. Oral and Maxillofacial Pathology (5th Ed).
% 3. PMID: 41841403 - Nguyen K et al. (Oral Dis 2026) Management and Outcomes of Recurrent Ranulas: Meta-Analysis.
% 4. PMID: 40921455 - Souza MRF et al. (J Oral Pathol Med 2025) Treatments to Avoid Ranula Recurrence: Network Meta-Analysis.
% 5. PMID: 42486748 - Barnes PM et al. (Br J Oral Maxillofac Surg 2026) Submandibular Duct Sialendoscopy Interventions.
% 6. PMID: 40087169 - Özçelik N et al. (Eur Arch Otorhinolaryngol 2025) Submandibular Sialolithiasis Holmium:YAG Lithotripsy.
% 7. PMID: 41479499 - Ayoub M, Ghosein E. (Cureus 2025) Management of Recurrent Lip Mucocele & Feeder Gland Excision.
% =========================================================================

\documentclass[a4paper]{article}
\usepackage[margin=1.15cm,footskip=0pt,headsep=0pt]{geometry}
\usepackage{amsmath,amssymb}
\usepackage{xcolor}
\usepackage{tikz}
\usetikzlibrary{arrows.meta,positioning,shapes.geometric}
\usepackage{hyperref}
\usepackage{fontspec}

\makeatletter
\newcommand{\MathCollectFalse}{}
\newcommand{\MathCollectTrue}{}
\let\@kernel@tabular@init\relax
\makeatother

\usepackage{xepersian}

\settextfont[
  Path = Peyda/,
  BoldFont = Peyda-Bold.ttf,
  ItalicFont = Peyda-Medium.ttf,
  BoldItalicFont = Peyda-Bold.ttf
]{Peyda-Regular.ttf}

\setlatintextfont{Times New Roman}

\IfFileExists{Peyda/Symbola.ttf}{%
  \newfontfamily\symb[Path=Peyda/]{Symbola.ttf}%
}{%
  \IfFileExists{editor/Peyda/Symbola.ttf}{%
    \newfontfamily\symb[Path=editor/Peyda/]{Symbola.ttf}%
  }{%
    \newfontfamily\symb{Symbola}%
  }%
}
\newcommand{\emoji}[1]{{\symb #1}}

\linespread{1.10}
\pagestyle{empty}
\thispagestyle{empty}

\raggedbottom
\setlength{\parindent}{0pt}
\setlength{\parskip}{0pt}

\definecolor{hlgray}{HTML}{E2E8F0}
\definecolor{hlyellow}{HTML}{FEF08A}
\definecolor{hlred}{HTML}{FEE2E2}
\definecolor{darkred}{HTML}{991B1B}
\definecolor{navy}{HTML}{0A2540}
\definecolor{teal}{HTML}{0D9488}

\begin{document}

\begin{center}
{\Large\bfseries فلوچارت تصمیم‌گیری بالینی: ضایعات تروماتیک و انسدادی بزاقی دهان}\\[2pt]
{\normalsize\bfseries\color{navy} \lr{Traumatic and Obstructive Salivary Lesions of the Oral Cavity (Differential Triage)}}\\[3pt]
{\small\itshape تریاژ موکوسل و غدد فرعی، رانولای کف دهان، سیالولیتیازیس مجرا (\lr{Sialendoscopy}) و نوروم تروماتیک دردناک}
\end{center}
\vspace{-4pt}
\hrule height 0.8pt
\vspace{6pt}

\begin{center}
\begin{latin}
\linespread{0.95}\selectfont
\begin{tikzpicture}[
    every node/.style={font=\sffamily\footnotesize, align=center},
    stepbox/.style={rectangle, draw=black, fill=white, line width=0.7pt, inner sep=4.5pt},
    decision/.style={diamond, aspect=2.4, draw=black, fill=white, line width=0.7pt, inner sep=2pt, font=\sffamily\scriptsize\bfseries},
    arrow/.style={-{Stealth[scale=0.9]}, thick, black}
]

% 1. Initial Presentation
\node (n1) [stepbox, line width=1.0pt, text width=13.0cm] at (0, 0.0) {
    \rl{\textbf{تظاهر بالینی: تورم کیستیک با نوسان اندازه حین غذا خوردن، توده زیرمخاطی دردناک یا سنگ مجرا}}\\[2pt]
    \scriptsize \rl{تاول شفاف متمایل به آبی در لب پایین یا کف دهان $\bullet$ درد و اتساع متناوب غدد ساب‌مندیبولار/پاروتید $\bullet$ سابقه تروما/جراحی}
};

% 2. Mandatory Step 1
\node (n2) [stepbox, text width=14.2cm] at (0, -1.60) {
    \rl{\colorbox{hlyellow}{\strut\textbf{مرحله ۱ الزامی: ارزیابی ترشح بزاق با دوشش مجرا (\lr{Milking})، رادیوگرافی اکلوزال و سونوگرافی}}}\\[2pt]
    \scriptsize \rl{بررسی جریان شفاف یا چرکی بزاق از سوراخ وارتون/استنسن $\bullet$ کشف کلسیفیکاسیون رادیواپک یا اتساع داکت در سونوگرافی}
};

% 3. Solitary Salivary Infection Decision
\node (d1) [decision] at (0, -3.30) {\rl{ترشح چرکی یا سلولیت حاد؟}\\[-1pt]\scriptsize \lr{Acute Infection?}};
\node (r_benign) [stepbox, draw=teal, line width=0.9pt, text width=3.8cm] at (5.8, -3.30) {
    \rl{\textbf{سیالادنیت حاد باکتریال \emoji{🔍}}}\\[2pt]
    \scriptsize \rl{آنتی‌بیوتیک سیستمیک؛}\\\rl{هیدراتاسیون و سیالوگوگ؛}\\\rl{پرهیز از پروپ‌زدن در فاز حاد}
};

% 4. Mandatory Mapping
\node (n3) [stepbox, line width=1.0pt, text width=15.8cm] at (0, -5.00) {
    \rl{\colorbox{hlyellow}{\strut\textbf{تریاژ موضع آناتومیک، عمق گرفتگی بافت نرم و تمایز سنگ داکت از اکستراوازیشن موکوس \emoji{🔬}}}}\\[2pt]
    \scriptsize \rl{بررسی گسترش رانولا به زیر عضله میلوهیوئید (\lr{Plunging Ranula}) $\bullet$ ارزیابی اندازه سنگ جهت انتخاب تکنیک آندوسکوپی \lr{[\textbf{PMID: 40087169}]}}
};

% 5. Subtype Decision
\node (d2) [decision] at (0, -6.60) {\rl{فنوتیپ بالینی}\\[-1pt]\scriptsize \lr{Clinical Entity?}};

% 6. 4 Main Columns
\node (h1) [stepbox, text width=3.6cm] at (-6.15, -9.00) {
    \rl{\textbf{موکوسل دهانی}}\\[2pt]
    \scriptsize \lr{Oral Mucocele (Extravasation)}\\\rl{پارگی تروماتیک داکت غدد فرعی؛}\\\rl{تجمع موین بدون پوشش اپیتلیال؛}\\\rl{عود با آسپیراسیون:} \colorbox{hlgray}{\strut\textbf{\rl{نزدیک ۱۰۰٪}}}
};
\node (h2) [stepbox, draw=darkred, line width=1.2pt, text width=3.6cm] at (-2.05, -9.00) {
    \rl{\textbf{رانولا (کف دهان / گردنی) \emoji{⚡}}}\\[2pt]
    \scriptsize \lr{Ranula / Plunging Ranula}\\\rl{نشت موکوس از غده زیرزبانی؛}\\\rl{اتساع یکطرفه حجیم زیر زبان؛}\\\rl{عود با مارسوپیالیزاسیون:} \colorbox{hlred}{\strut\color{darkred}\textbf{\rl{۳۰٪ تا ۵۰٪}}}
};
\node (h3) [stepbox, text width=3.6cm] at (2.05, -9.00) {
    \rl{\textbf{سیالولیتیازیس (سنگ مجرا)}}\\[2pt]
    \scriptsize \lr{Salivary Sialolithiasis}\\\rl{رسوب نمک‌های کلسیم در مجرای وارتون؛}\\\rl{اتساع دردناک غده حین صرف غذا؛}\\\rl{موفقیت سیالندوسکوپی: \colorbox{hlgray}{\strut\textbf{\rl{بالای ۸۵٪}}}}
};
\node (h4) [stepbox, text width=3.6cm] at (6.15, -9.00) {
    \rl{\textbf{نوروم تروماتیک فکین}}\\[2pt]
    \scriptsize \lr{Traumatic Neuroma}\\\rl{پرولیفراسیون اکسون ناشی از تروما؛}\\\rl{ندول دردناک در فورامن منتال/لب؛}\\\rl{علامت مثبت تینل (\lr{Tinel Sign})}
};

% 7. Primary Interventions
\node (i1) [stepbox, text width=3.6cm] at (-6.15, -11.60) {
    \rl{\textbf{اکسیزیون موکوسل + غدد فرعی}}\\[2pt]
    \scriptsize \rl{برداشت کامل لومن سودوسیست؛}\\\rl{تراش و خارج‌سازی غدد بزاقی مجاور؛}\\\rl{کاهش عود به زیر ۲٪}\\\rl{\lr{[\textbf{PMID: 41479499}]}}
};
\node (i2) [stepbox, draw=darkred, line width=1.2pt, text width=3.6cm] at (-2.05, -11.60) {
    \rl{\textbf{خارج‌سازی غده زیرزبانی}}\\[2pt]
    \scriptsize \textbf{\rl{رزکسیون ساب‌لینگوال اکلکتومی؛}}\\\rl{حفظ عصب لینگوال و مجرای وارتون؛}\\\rl{درمان قطعی استاندارد رانولا}\\\rl{\lr{[\textbf{PMID: 41841403}]}}
};
\node (i3) [stepbox, text width=3.6cm] at (2.05, -11.60) {
    \rl{\textbf{سیالندوسکوپی $\pm$ لیتوتریپسی}}\\[2pt]
    \scriptsize \rl{سنگ‌های زیر ۵mm: خارج‌سازی با بسکت؛}\\\rl{سنگ‌های بالای ۵mm: خرد کردن با لیزر،}\\\rl{یا داکتوتومی ترانس‌اورال در کف دهان}\\\rl{\lr{[\textbf{PMID: 42486748}]}}
};
\node (i4) [stepbox, text width=3.6cm] at (6.15, -11.60) {
    \rl{\textbf{اکسیزیون جراحی تروماتیک}}\\[2pt]
    \scriptsize \rl{برداشت ندول همراه با حاشیه کوچک عصب؛}\\\rl{کوتر پروکسیمال تنه عصب جهت}\\\rl{پیشگیری از تشکیل مجدد جوانه نوروما}
};

% 8. Surveillance Regimens
\node (s1) [stepbox, text width=3.6cm] at (-6.15, -14.20) {
    \rl{\textbf{زمان‌بندی پیگیری:}}\\[2pt]
    \rl{\textbf{هر ۶ ماه تا ۱ سال}}\\[2pt]
    \scriptsize \rl{بررسی مخاط لب پایین؛}\\\rl{عدم بروز بی‌حسی لب یا عود موکوسل}
};
\node (s2) [stepbox, draw=darkred, line width=1.2pt, text width=3.6cm] at (-2.05, -14.20) {
    \rl{\textbf{زمان‌بندی پیگیری:}}\\[2pt]
    \colorbox{hlred}{\strut\color{darkred}\textbf{\rl{هر ۳ ماه تا ۲ سال}}}\\[2pt]
    \scriptsize \rl{معاینه دقیق کف دهان و گردن؛}\\\rl{سونوگرافی بافت نرم در رانولای پلانجینگ}
};
\node (s3) [stepbox, text width=3.6cm] at (2.05, -14.20) {
    \rl{\textbf{زمان‌بندی پیگیری:}}\\[2pt]
    \rl{\textbf{هر ۶ ماه تا ۲ سال}}\\[2pt]
    \scriptsize \rl{ارزیابی پتنسی داکت بزاقی؛}\\\rl{مصرف مایعات و ماساژ منظم غده}
};
\node (s4) [stepbox, text width=3.6cm] at (6.15, -14.20) {
    \rl{\textbf{زمان‌بندی پیگیری:}}\\[2pt]
    \rl{\textbf{سالانه تا ۲ سال}}\\[2pt]
    \scriptsize \rl{پایش دردهای نوروپاتیک؛}\\\rl{درمان تکمیلی دارویی در صورت نیاز}
};

% 9. Clinical Red Flags
\node (redflags) [stepbox, draw=darkred, line width=0.9pt, text width=15.8cm] at (0, -17.40) {
    \rl{\textbf{هشدارهای بالینی و اندیکاسیون‌های ارجاع فوری یا جراحی اورژانسی (\lr{Clinical Red Flags}):}}\\[3pt]
    \scriptsize \colorbox{hlred}{\strut\color{darkred}\textbf{\rl{مداخله فوری در:}}} \rl{(۱) اتساع رانولای پلانجینگ به فضای پارافارنژیال و بروز تنگی نفس یا دیسفاژی شدید $\bullet$ (۲) سیالادنیت چرکی حاد پیشرونده با تریسموس و درگیری فضاهای عمقی گردن (آنژین لودویگ) $\bullet$ (۳) آسیب حین عمل به شاخه مارجینال مندیبولار عصب فاشیال یا عصب زبانی $\bullet$ (۴) پایداری توده کف دهان بدون ارتباط با غذا خوردن مشکوک به موکواپیدرموئید کارسینوما.}
};

% 10. Refractory / Malignancy Protocols
\node (refractory) [stepbox, line width=0.9pt, text width=15.8cm] at (0, -20.00) {
    \rl{\textbf{پروتکل موارد راجعه و درمان‌های پیشرفته بزاقی (\lr{Advanced Salivary Protocols}):}}\\[3pt]
    \scriptsize \rl{\textbf{سطح ۱ (رانولای راجعه پس از مارسوپیالیزاسیون):} خارج‌سازی کامل و بی قید و شرط غده زیرزبانی در اتاق عمل همراه با درناژ کامل حفرات فاسیال گردن \lr{[\textbf{PMID: 40921455}]}.}\\
    \scriptsize \rl{\textbf{سطح ۲ (سیالولیتیازیس پارانشیمال غده ساب‌مندیبولار):} در موارد سنگ‌های متعدد در عمق پارانشیم با آتروفی فیبروتیک غده، سیالادنکتومی کامل با دسترسی جراحی زیر فک.}\\
    \scriptsize \rl{\textbf{سطح ۳ (نوروپاتی راجعه و دردهای مزمن پس از قطع عصب):} کاربرد تکنیک‌های بازسازی میکروجراحی عصب (\lr{Nerve Graft / Conduit}) یا درمان دارویی با پره‌گابالین و دوپوکستین.}
};

% Connections
\draw [arrow] (n1) -- (n2);
\draw [arrow] (n2) -- (d1);
\draw [arrow] (d1.east) -- node[above, font=\scriptsize\bfseries] {\rl{خاموش / خوش‌خیم}} (r_benign.west);
\draw [arrow] (d1.south) -- node[right=5pt, font=\scriptsize\bfseries] {\rl{فعال / مشکوک}} (n3.north);
\draw [arrow] (n3.south) -- (d2.north);

% Bus-Bar from Decision 2 to 4 Subtype columns
\draw [thick, black] (d2.south) -- (0, -7.50);
\draw [thick, black] (-6.15, -7.50) -- (6.15, -7.50);
\draw [arrow] (-6.15, -7.50) -- (h1.north);
\draw [arrow] (-2.05, -7.50) -- (h2.north);
\draw [arrow] (2.05, -7.50) -- (h3.north);
\draw [arrow] (6.15, -7.50) -- (h4.north);

% Column Downward Arrows
\draw [arrow] (h1.south) -- (i1.north);
\draw [arrow] (h2.south) -- (i2.north);
\draw [arrow] (h3.south) -- (i3.north);
\draw [arrow] (h4.south) -- (i4.north);

\draw [arrow] (i1.south) -- (s1.north);
\draw [arrow] (i2.south) -- (s2.north);
\draw [arrow] (i3.south) -- (s3.north);
\draw [arrow] (i4.south) -- (s4.north);

% Collector Bus into Trigger Box
\draw [thick, black] (s1.south) -- (-6.15, -15.80);
\draw [thick, black] (s2.south) -- (-2.05, -15.80);
\draw [thick, black] (s3.south) -- (2.05, -15.80);
\draw [thick, black] (s4.south) -- (6.15, -15.80);
\draw [thick, black] (-6.15, -15.80) -- (6.15, -15.80);
\draw [arrow] (0, -15.80) -- (redflags.north);

% Final link to Refractory Box
\draw [arrow] (redflags.south) -- (refractory.north);
\end{tikzpicture}
\end{latin}
\end{center}

\end{document}
